Learning Reflection: This was another bachelor degree's worth of learning in itself. On the skills side of things, we learned how to practically apply kinematics, reinforcement learning, robot perception, environment simulation, and so many more topics to a real, tangible problem. Aside from that, we had to familiarize ourselves extensively with novel softwares and their quirks particularly ROS2 and all encompassed under it. Of course, it wasn't all hard skills; we learned how to cleanly document our work and regularly present our findings per week in a clear and confident manner.

Team Reflection: Despite having an interdisciplinary team, we organized our efforts based on our strengths. Moreover, we made sure to maintain a clear, shared vision by checking in on each other in times of stress while ensuring we were all agreed on doing what we wanted to do before moving forward. All this proved indispensable in propelling our project forward.

Project Reflection: It is great that our project has fulfilled the individual segments that constitute our deliverables. We had planned on extending perception into the real domain, but due to lack of time as well as poor motor durability, that could not happen. However, on the flip side, there were multiple ambitious parts particularly reinforcement learning we never thought we would be able to accomplish, but we did.

Process Reflection: We were efficient and our task divisions normally clear. To make sure everyone was on the same page, we maintained a OneNote document where each teammate documented whatever work they have done, what work they will do, and any current or potential obstacles. What could have gone better was to concentrate efforts in first accomplishing the general robot pipeline before trying out various approaches for smaller segments.
